% Part: first-order-logic
% Chapter: syntax-and-semantics
% Section: Structures

\documentclass[../../../include/open-logic-section]{subfiles}

\begin{document}

\olfileid{fol}{syn}{str}

\olsection{\printtoken{P}{structure} للغات الرتبة الأولى}

\begin{explain}
لغة الرتبة الأولى إذا نُظر إليها في ذاتها كانت \emph{غير مؤوَّلة}؛
إذ لا معنى معيّنًا لرموز ال!!{constant}s وال!!{function}s
وال!!{predicate}s بمجردها. وإنما نعيّن معانيها باختيار
\article{structure} \emph{!!{structure}}. ففيها يتعيّن
\emph{المجال}: الأشياء التي تدل عليها ال!!{constant}s، وتجري
عليها ال!!{function}s، وتتراوح عليها المكممات. ويتعيّن فيها كذلك
مدلول كل واحد من ال!!{constant}s، وكيف تنقل !!a{function} الأشياء
إلى أشياء، وعلى أي الأشياء تصدق ال!!{predicate}s.
وعلى ال!!^{structure}s تُبنى المفاهيم \emph{الدلالية} في المنطق،
كالنتيجة والصلاحية وقابلية الإرضاء. وقد تختلف تسميتها في
المؤلفات، فيقال «بنى» أو «تأويلات» أو «نماذج».
\end{explain}

\begin{defn}[!!^{structure}s]
تتألف \Article{structure} \emph{!!{structure}}~$\Struct M$ للغة
$\Lang{L}$ من منطق الرتبة الأولى من العناصر الآتية:
\begin{enumerate}
\item \emph{المجال:} مجموعة غير خالية، $\Domain {\OLId{latin-looped}{M}}$
\item \emph{تأويل ال!!{constant}s:} لكل !!{constant}~${\OLId{latin-ordinary}{c}}$ من
  $\Lang{L}$، !!a{element}~$\Assign{{\OLId{latin-ordinary}{c}}}{{\OLId{latin-looped}{M}}} \in \Domain {\OLId{latin-looped}{M}}$
\item \emph{تأويل ال!!{predicate}s:} لكل !!{predicate} ذات
  ${\OLId{latin-ordinary}{n}}$ مواضع~${\OLId{latin-looped}{R}}$ من~$\Lang{L}$ (عدا~$\eq$)، علاقة ذات ${\OLId{latin-ordinary}{n}}$ مواضع
  $\Assign{{\OLId{latin-looped}{R}}}{{\OLId{latin-looped}{M}}} \subseteq \Domain{{\OLId{latin-looped}{M}}}^{\OLId{latin-ordinary}{n}}$
\item \emph{تأويل ال!!{function}s:} لكل !!{function} ذات
  ${\OLId{latin-ordinary}{n}}$ مواضع~${\OLId{latin-ordinary}{f}}$ من~$\Lang{L}$، دالة ذات ${\OLId{latin-ordinary}{n}}$ مواضع
  $\Assign{{\OLId{latin-ordinary}{f}}}{{\OLId{latin-looped}{M}}} \colon \Domain{{\OLId{latin-looped}{M}}}^{\OLId{latin-ordinary}{n}} \to \Domain{{\OLId{latin-looped}{M}}}$
\end{enumerate}
\end{defn}

\begin{ex}
تتألف !!^a{structure}~$\Struct M$ للغة الحساب من مجموعة، ومن عنصر
من~$\Domain {\OLId{latin-looped}{M}}$، هو $\Assign{\Obj 0}{{\OLId{latin-looped}{M}}}$، تأويلًا لل!!{constant}
~$\Obj 0$، ومن دالة أحادية المحل
$\Assign{\Obj \prime}{{\OLId{latin-looped}{M}}} \colon \Domain{{\OLId{latin-looped}{M}}} \to \Domain {\OLId{latin-looped}{M}}$،
ودالتين ثنائيتي المحل $\Assign{\Obj +}{{\OLId{latin-looped}{M}}}$ و$\Assign{\Obj
  \times}{{\OLId{latin-looped}{M}}}$، كلتاهما $\Domain {\OLId{latin-looped}{M}}^{\OLMathNumeral{2}} \to \Domain {\OLId{latin-looped}{M}}$، وعلاقة ثنائية
المحل $\Assign{\Obj <}{{\OLId{latin-looped}{M}}} \subseteq \Domain{{\OLId{latin-looped}{M}}}^{\OLMathNumeral{2}}$.

وأوضح مثال لهذه البنية أن نختار ما يأتي:
\begin{enumerate}
\item $\Domain {\OLId{latin-looped}{N}} = \Nat$
\item $\Assign{\Obj 0}{{\OLId{latin-looped}{N}}} = {\OLMathNumeral{0}}$
\item $\Assign{\Obj \prime}{{\OLId{latin-looped}{N}}}({\OLId{latin-ordinary}{n}}) = {\OLId{latin-ordinary}{n}} + {\OLMathNumeral{1}}$ لكل ${\OLId{latin-ordinary}{n}} \in \Nat$
\item $\Assign{\Obj +}{{\OLId{latin-looped}{N}}}({\OLId{latin-ordinary}{n}}, {\OLId{latin-ordinary}{m}}) = {\OLId{latin-ordinary}{n}} + {\OLId{latin-ordinary}{m}}$ لكل ${\OLId{latin-ordinary}{n}}, {\OLId{latin-ordinary}{m}} \in \Nat$
\item $\Assign{\Obj \times}{{\OLId{latin-looped}{N}}}({\OLId{latin-ordinary}{n}}, {\OLId{latin-ordinary}{m}}) = {\OLId{latin-ordinary}{n}}\cdot {\OLId{latin-ordinary}{m}}$ لكل ${\OLId{latin-ordinary}{n}}, {\OLId{latin-ordinary}{m}} \in \Nat$
\item $\Assign{\Obj <}{{\OLId{latin-looped}{N}}} = \Setabs{\tuple{{\OLId{latin-ordinary}{n}}, {\OLId{latin-ordinary}{m}}}}{{\OLId{latin-ordinary}{n}} \in \Nat, {\OLId{latin-ordinary}{m}} \in
  \Nat, {\OLId{latin-ordinary}{n}} < {\OLId{latin-ordinary}{m}}}$
\end{enumerate}
وهذه البنية~$\Struct N$ للغة~$\Lang L_{\OLId{latin-looped}{A}}$ تسمى
\emph{النموذج القياسي للحساب}؛ فإن الرموز غير المنطقية
لـ~$\Lang L_{\OLId{latin-looped}{A}}$ تأخذ فيها المعاني الحسابية المعهودة.

وليست هذه كل ال!!{structure}s الممكنة لـ~$\Lang L_{\OLId{latin-looped}{A}}$، بل غيرها كثير.
فلنا أن نجعل الأعداد الصحيحة~$\Int$ مجالًا مكان~$\Nat$، ثم نختار
تأويلات~$\Obj 0$، و$\Obj \prime$، و$\Obj +$، و$\Obj
\times$، و$\Obj <$ على وفق ذلك. ولنا أيضًا أن نضع
لـ~$\Lang L_{\OLId{latin-looped}{A}}$ بنى لا صلة لها بالأعداد أصلًا.
\end{ex}

\begin{ex}
أما البنية~$\Struct M$ للغة نظرية المجموعات~$\Lang L_{\OLId{latin-looped}{Z}}$، فيكفي
لتعيينها مجموعة غير خالية وعلاقة ثنائية المحل. فيجوز، بمقتضى
التعريف الصوري، أن نأخذ مجموعة الناس مع العلاقة «${\OLId{latin-ordinary}{x}}$ أكبر سنًا من~${\OLId{latin-ordinary}{y}}$»
بوصفهما !!a{structure} لـ~$\Lang L_{\OLId{latin-looped}{Z}}$، كما يمكن استخدام~$\Nat$ مع
${\OLId{latin-ordinary}{n}} \ge {\OLId{latin-ordinary}{m}}$، حيث ${\OLId{latin-ordinary}{n}}, {\OLId{latin-ordinary}{m}} \in \Nat$.

ومن ال!!{structure}s الجديرة بالنظر للغة~$\Lang L_{\OLId{latin-looped}{Z}}$ ما تكون
!!{element}s مجاله مجموعات بالفعل، ويكون مدلول~$\Obj \in$
فيه نفس العلاقة «${\OLId{latin-ordinary}{x}}$ !!a{element} من~${\OLId{latin-ordinary}{y}}$»، ال!!{structure}
~$\Struct{{HF}}$ المؤلفة من \emph{المجموعات المنتهية وراثيًا}:
\begin{enumerate}
\item $\Domain{{{HF}}} = \emptyset \cup \Pow{\emptyset} \cup
  \Pow{\Pow{\emptyset}} \cup \Pow{\Pow{\Pow{\emptyset}}} \cup \dots$;
\item $\Assign{\Obj \in}{{{HF}}} = \Setabs{\tuple{{\OLId{latin-ordinary}{x}}, {\OLId{latin-ordinary}{y}}}}{{\OLId{latin-ordinary}{x}}, {\OLId{latin-ordinary}{y}} \in
  \Domain{{{HF}}}, {\OLId{latin-ordinary}{x}} \in {\OLId{latin-ordinary}{y}}}$.
\end{enumerate}
\end{ex}

\begin{digress}
وللشروط التي نضعها في تعريف !!a{structure} أثر في المنطق نفسه.
فمنع المجال الخالي يقتضي، مع المعنى المعتاد لإرضاء الجمل المكممة
أو صدقها، أن تكون
$\lexists[{\OLId{latin-ordinary}{x}}][(!A({\OLId{latin-ordinary}{x}}) \lor \lnot !A({\OLId{latin-ordinary}{x}}))]$ صالحة، أي حقيقة منطقية.
واشتراط أن تدل ال!!{constant}s كلها على أشياء من المجال يجعل
التعميم الوجودي استدلالًا سليمًا: من $!A({\OLId{latin-ordinary}{a}})$ يلزم
$\lexists[{\OLId{latin-ordinary}{x}}][!A({\OLId{latin-ordinary}{x}})]$. أما إن أجزنا أن تدل الأسماء على ما خارج المجال،
أو ألا تدل على شيء، فقد أخذنا بمسلك \emph{المنطق الحر}؛
وفيه يحتاج التعميم الوجودي إلى مقدمة زائدة: من $!A({\OLId{latin-ordinary}{a}})$
و$\lexists[{\OLId{latin-ordinary}{x}}][\eq[{\OLId{latin-ordinary}{x}}][{\OLId{latin-ordinary}{a}}]]$ يلزم
$\lexists[{\OLId{latin-ordinary}{x}}][!A({\OLId{latin-ordinary}{x}})]$.
\end{digress}

\end{document}
